% =========================================================
% Lattices and Complete Lattices
% =========================================================

\section{Lattices}\label{sec:lattices}

% ---------------------------------------------------------
% TOOLKIT
% ---------------------------------------------------------
\begin{tcolorbox}[colback=gray!6, colframe=gray!40, arc=2pt,
  left=6pt, right=6pt, top=4pt, bottom=4pt,
  title={\small\textbf{Lattices — Quick Reference}},
  fonttitle=\small\bfseries]
\small
\begin{tabular}{l l l}
\toprule
\textbf{Concept} & \textbf{Meaning} & \textbf{Detail} \\
\midrule
Lattice              & Poset with $\sup\{a,b\}$ and $\inf\{a,b\}$ for all $a,b$   & \hyperref[def:lattice]{↓ Def} \\
Join / Meet          & $a \vee b = \sup\{a,b\}$; $a \wedge b = \inf\{a,b\}$        & \hyperref[def:lattice]{↓ Def} \\
Complete lattice     & Every subset has sup and inf                                & \hyperref[def:complete-lattice]{↓ Def} \\
Distributive lattice & $a \wedge (b \vee c) = (a \wedge b) \vee (a \wedge c)$      & \hyperref[def:distributive-lattice]{↓ Def} \\
Boolean lattice      & Complemented distributive lattice                           & \hyperref[def:boolean-lattice]{↓ Def} \\
$\mathcal{P}(X)$     & Complete Boolean lattice under $\subseteq$                  & \hyperref[ex:power-set-lattice]{↓ Ex} \\
$\sigma$-algebras    & Sub-lattices of $\mathcal{P}(X)$ closed under countable ops  & \hyperref[rem:sigma-alg-lattice]{↓ Rem} \\
Knaster--Tarski      & Every monotone map on a complete lattice has a fixed point   & \hyperref[thm:knaster-tarski]{↓ Thm} \\
\bottomrule
\end{tabular}
\end{tcolorbox}

\vspace{1em}

% ---------------------------------------------------------
% Motivation
% ---------------------------------------------------------

The preceding section defined supremum and infimum for arbitrary subsets of a
poset.  In many posets that appear in practice, every pair of elements already
has a supremum and an infimum.  When this is the case, the binary operations
$a \vee b = \sup\{a,b\}$ and $a \wedge b = \inf\{a,b\}$ are well-defined
everywhere, giving the poset the algebraic structure of a \emph{lattice}.

Lattices are ubiquitous:
\begin{itemize}
  \item The power set $\mathcal{P}(X)$ ordered by $\subseteq$ is a complete
        lattice, with join = union and meet = intersection.
  \item Every $\sigma$-algebra is a sub-lattice of $\mathcal{P}(X)$
        that is additionally closed under countable joins and meets and under
        complementation.
  \item Every topology is a sub-structure of $\mathcal{P}(X)$ that is closed
        under arbitrary joins and finite meets.
  \item In functional analysis, the lattice of closed subspaces of a Hilbert
        space is a complete lattice under inclusion.
\end{itemize}
Understanding lattices here means that when $\sigma$-algebras, topologies, and
sub-$\sigma$-algebras are encountered downstream, their closure conditions
will be recognisable instances of a single abstract pattern.

% ---------------------------------------------------------
% Lattice definition
% ---------------------------------------------------------

\begin{tcolorbox}[colback=propbox, colframe=propborder, arc=2pt,
  left=6pt, right=6pt, top=4pt, bottom=4pt,
  title={\small\textbf{Definition (Lattice)}},
  fonttitle=\small\bfseries]
\label{def:lattice}
A \emph{lattice} is a poset $(L, \leq)$ in which every two-element subset
$\{a, b\} \subseteq L$ has both a supremum and an infimum in $L$.

The supremum $\sup\{a,b\}$ is called the \emph{join} of $a$ and $b$,
written $a \vee b$.
The infimum $\inf\{a,b\}$ is called the \emph{meet} of $a$ and $b$,
written $a \wedge b$.
\end{tcolorbox}

\begin{remark*}[Fully quantified form]
$(L, \leq)$ is a lattice $\;\leftrightarrow\;$
$\forall a, b \in L,\;\exists\, a \vee b \in L \;\text{and}\; \exists\, a \wedge b \in L$,
where $a \vee b$ and $a \wedge b$ satisfy the sup and inf conditions of
\hyperref[def:sup-inf]{Definition~(Supremum and Infimum)}.
\end{remark*}

\begin{remark}[Join and meet as operations]
\label{rem:lattice-operations}
In a lattice, $\vee$ and $\wedge$ are binary operations $L \times L \to L$.
They satisfy, for all $a, b, c \in L$:
\begin{enumerate}[label=(\roman*)]
  \item \emph{Commutativity}: $a \vee b = b \vee a$ and $a \wedge b = b \wedge a$.
  \item \emph{Associativity}: $(a \vee b) \vee c = a \vee (b \vee c)$ and
        $(a \wedge b) \wedge c = a \wedge (b \wedge c)$.
  \item \emph{Idempotency}: $a \vee a = a$ and $a \wedge a = a$.
  \item \emph{Absorption}: $a \vee (a \wedge b) = a$ and
        $a \wedge (a \vee b) = a$.
\end{enumerate}
Properties (i)--(iv) follow directly from the order-theoretic definitions of join
and meet.  Conversely, any set with two binary operations satisfying (i)--(iv)
can be given a partial order making it a lattice: define $a \leq b$ iff
$a \wedge b = a$.  The two approaches — order-theoretic and algebraic — are
equivalent.
\end{remark}

\begin{example}[Power set lattice]
\label{ex:power-set-lattice}
For any set $X$, the power set $(\mathcal{P}(X), \subseteq)$ is a lattice:
$A \vee B = A \cup B$ and $A \wedge B = A \cap B$.  The algebraic laws
of the Set Algebra section (commutativity, associativity, distributivity,
absorption) are exactly the lattice laws in this example.  This is not a
coincidence: Boolean algebra is a special kind of lattice.
\end{example}

\begin{example}[Divisibility lattice]
\label{ex:div-lattice}
The divisibility order on $\mathbb{N}^+ = \{1, 2, 3, \ldots\}$ makes it a
lattice: $a \vee b = \mathrm{lcm}(a,b)$ and $a \wedge b = \gcd(a,b)$, as
we noted in Example~\ref{ex:sup-divisibility}.
\end{example}

\begin{example}[Not every poset is a lattice]
\label{ex:non-lattice}
Consider the poset $(\{a,b,c,d\}, \leq)$ with $a < c$, $a < d$, $b < c$,
$b < d$, and $a,b$ incomparable.  The subset $\{a,b\}$ has two minimal
upper bounds, $c$ and $d$, which are incomparable.  Hence $\sup\{a,b\}$
does not exist, and this poset is not a lattice.

The power set $\mathcal{P}(\{1,2,3\})$ restricted to subsets of cardinality
$\neq 1$ provides a cleaner example: $\{1,2\}$ and $\{1,3\}$ have no meet in
this restricted collection.
\end{example}

% ---------------------------------------------------------
% Complete Lattice
% ---------------------------------------------------------

\begin{tcolorbox}[colback=propbox, colframe=propborder, arc=2pt,
  left=6pt, right=6pt, top=4pt, bottom=4pt,
  title={\small\textbf{Definition (Complete Lattice)}},
  fonttitle=\small\bfseries]
\label{def:complete-lattice}
A lattice $(L, \leq)$ is a \emph{complete lattice} if every subset
$S \subseteq L$ (including the empty set and $L$ itself) has both a
supremum and an infimum in $L$:
\[
  \forall S \subseteq L,\quad \sup S \in L \;\text{ and }\; \inf S \in L.
\]
\end{tcolorbox}

\begin{remark*}[Fully quantified form]
$(L, \leq)$ is a complete lattice $\;\leftrightarrow\;$
$\forall S \subseteq L,\;\exists u^* \in L,\;\exists \ell^* \in L$ satisfying
the sup and inf conditions of \hyperref[def:sup-inf]{Definition~(Supremum and Infimum)}.
\end{remark*}

\begin{remark}[The empty set forces top and bottom]
\label{rem:complete-top-bottom}
When $S = \varnothing$, the sup condition says: $u^*$ is an upper bound of
the empty set (which is vacuously satisfied by every element of $L$), and $u^*$
is at most every other upper bound.  This forces $u^* = \inf L$, the
\emph{bottom element} of $L$ (often written $\bot$).  Symmetrically,
$\inf \varnothing = \sup L$, the \emph{top element} (often written $\top$).

Concretely: in $(\mathcal{P}(X), \subseteq)$, we have $\sup \varnothing =
\varnothing$ (the empty union) and $\inf \varnothing = X$ (the empty
intersection equals the whole universe, by convention).  Every complete
lattice therefore has a top element $\top = \sup L$ and a bottom element
$\bot = \inf L$.
\end{remark}

\begin{example}[$\mathcal{P}(X)$ is a complete lattice]
\label{ex:power-set-complete}
For any set $X$, the power set $(\mathcal{P}(X), \subseteq)$ is a complete
lattice: for any family $\{A_\alpha\}_{\alpha \in I} \subseteq \mathcal{P}(X)$,
\[
  \sup_{\alpha \in I} A_\alpha = \bigcup_{\alpha \in I} A_\alpha,
  \qquad
  \inf_{\alpha \in I} A_\alpha = \bigcap_{\alpha \in I} A_\alpha.
\]
The bottom element is $\varnothing$ and the top element is $X$.  This is the
canonical example of a complete lattice, and every $\sigma$-algebra on $X$ is a
complete sub-lattice of $(\mathcal{P}(X), \subseteq)$.
\end{example}

\begin{example}[\texorpdfstring{$[0,\infty]$}{[0,inf]} is a complete lattice]
\label{ex:extended-positive}
The extended positive half-line $[0,\infty] = [0,\infty) \cup \{+\infty\}$
with the usual order is a complete lattice.  Every nonempty subset has a
supremum (which may be $+\infty$) and an infimum (which may be $0$).  This
is the value space for measures: a measure assigns to each measurable set
a value in $[0,\infty]$, and the key properties of measures (countable
additivity, monotone convergence) are stated in terms of sups and infs in
this complete lattice.
\end{example}

% ---------------------------------------------------------
% Distributive and Boolean lattices
% ---------------------------------------------------------

\begin{tcolorbox}[colback=propbox, colframe=propborder, arc=2pt,
  left=6pt, right=6pt, top=4pt, bottom=4pt,
  title={\small\textbf{Definition (Distributive Lattice)}},
  fonttitle=\small\bfseries]
\label{def:distributive-lattice}
A lattice $(L, \leq)$ is \emph{distributive} if for all $a, b, c \in L$:
\[
  a \wedge (b \vee c) = (a \wedge b) \vee (a \wedge c).
\]
\end{tcolorbox}

\begin{remark}[One law implies the other]
\label{rem:dist-dual}
By the duality principle for lattices, distributivity of $\wedge$ over $\vee$
is equivalent to distributivity of $\vee$ over $\wedge$:
$a \vee (b \wedge c) = (a \vee b) \wedge (a \vee c)$.
Proving one gives the other for free.
\end{remark}

\begin{remark}[Why distributivity matters]
\label{rem:dist-matters}
The distributive law is what makes a lattice behave like set algebra.  In the
power set $\mathcal{P}(X)$, the distributive laws are exactly the set
identities $A \cap (B \cup C) = (A \cap B) \cup (A \cap C)$ from
Theorem~3.22.  Distributivity is the condition that separates lattices that
``look like sets'' from those that do not (e.g., the lattice of closed
subspaces of a Hilbert space is not distributive).
\end{remark}

\begin{tcolorbox}[colback=propbox, colframe=propborder, arc=2pt,
  left=6pt, right=6pt, top=4pt, bottom=4pt,
  title={\small\textbf{Definition (Complement and Boolean Lattice)}},
  fonttitle=\small\bfseries]
\label{def:boolean-lattice}
Let $(L, \leq)$ be a bounded distributive lattice with top element $\top$ and
bottom element $\bot$.  An element $a' \in L$ is a \emph{complement} of
$a \in L$ if
\[
  a \vee a' = \top \quad \text{and} \quad a \wedge a' = \bot.
\]
A \emph{Boolean lattice} (or \emph{Boolean algebra}) is a bounded distributive
lattice in which every element has a complement.
\end{tcolorbox}

\begin{remark}[The power set is the canonical Boolean lattice]
\label{rem:power-set-boolean}
In $\mathcal{P}(X)$, the complement of a set $A$ is $A^c = X \setminus A$,
since $A \cup A^c = X = \top$ and $A \cap A^c = \varnothing = \bot$.  The
power set is therefore a complete Boolean lattice.  Every $\sigma$-algebra
on $X$ is a Boolean sub-algebra of $\mathcal{P}(X)$: it is closed under
complements (by definition) and under finite meets and joins (since every
$\sigma$-algebra is an algebra).
\end{remark}

% ---------------------------------------------------------
% Sub-lattices and sigma-algebras
% ---------------------------------------------------------

\begin{remark}[$\sigma$-algebras as sub-lattices]
\label{rem:sigma-alg-lattice}
A $\sigma$-algebra $\mathcal{M}$ on $X$ is, from the lattice-theoretic
viewpoint, a sub-structure of $(\mathcal{P}(X), \subseteq)$ that is:
\begin{itemize}
  \item a complete sub-lattice (closed under all countable joins and meets);
  \item a Boolean sub-algebra (closed under complementation).
\end{itemize}
The condition ``closed under countable joins'' is exactly the $\sigma$-algebra
axiom that countable unions of measurable sets are measurable.  The condition
``closed under meets'' for countable families follows from De Morgan together
with closure under countable unions and complements.

This perspective makes the definition of $\sigma$-algebra feel natural: it is
exactly a complete Boolean sub-algebra of the power set lattice, where
``complete'' is interpreted countably (countable joins and meets), not
set-theoretically.
\end{remark}

\begin{remark}[Topologies as join-sub-structures]
\label{rem:topology-lattice}
A topology $\tau$ on $X$ is \emph{not} a sub-lattice of $\mathcal{P}(X)$ in
general, because it is not closed under arbitrary meets (i.e., arbitrary
intersections of open sets need not be open).  However, $\tau$ is closed under
arbitrary joins (arbitrary unions of open sets are open) and finite meets
(finite intersections of open sets are open).  This makes $\tau$ an
\emph{inf-semilattice} for finite meets and a complete \emph{join-semilattice}.
The failure of closure under infinite meets is fundamental: it is precisely
what allows open sets to ``separate'' points, and it is the reason topology
is structurally different from measure theory.
\end{remark}

% ---------------------------------------------------------
% Knaster-Tarski
% ---------------------------------------------------------

\begin{tcolorbox}[colback=thmbox, colframe=thmborder, arc=2pt,
  left=6pt, right=6pt, top=4pt, bottom=4pt,
  title={\small\textbf{Theorem (Knaster--Tarski Fixed-Point Theorem)}},
  fonttitle=\small\bfseries]
\label{thm:knaster-tarski}
Let $(L, \leq)$ be a complete lattice and $f : L \to L$ a monotone
function (i.e., $a \leq b \Rightarrow f(a) \leq f(b)$).  Then $f$ has a
least fixed point and a greatest fixed point in $L$.  Specifically:
\[
  \mathrm{lfp}(f) = \inf\{\, a \in L : f(a) \leq a \,\},
  \qquad
  \mathrm{gfp}(f) = \sup\{\, a \in L : a \leq f(a) \,\}.
\]
\end{tcolorbox}

\begin{remark*}[Fully quantified form]
$\forall$ complete lattices $(L,\leq)$, $\forall$ monotone $f : L \to L$:
$\exists p, q \in L$ such that $f(p) = p$, $f(q) = q$, and $\forall r \in L$
with $f(r) = r$: $p \leq r \leq q$.
\end{remark*}

\begin{remark}[Why this theorem appears here]
\label{rem:knaster-tarski-context}
The Knaster--Tarski theorem is not needed for most of Volume~I, but it is
placed here because its hypotheses — complete lattice, monotone function —
are now in place and the statement is short.  Its significance reaches
further:
\begin{itemize}
  \item In set theory, it is used to prove the Schröder--Bernstein theorem
        without the Axiom of Choice (the monotone map sends $A \mapsto
        X \setminus g(Y \setminus f(A))$ on $\mathcal{P}(X)$).
  \item In measure theory, it is used to construct the smallest
        $\sigma$-algebra containing a given class of sets (the
        generated $\sigma$-algebra is a fixed point of the ``closure
        under countable operations'' operator).
  \item In theoretical computer science and domain theory, it is the
        foundation for denotational semantics.
\end{itemize}
For now, note that $(\mathcal{P}(X), \subseteq)$ is a complete lattice and
any set operation that is monotone (such as $A \mapsto \overline{A}$, closure
in a topological space, or $A \mapsto \sigma(A)$, generated $\sigma$-algebra)
has a fixed point by this theorem.
\end{remark}

\begin{remark}[Transition]
\label{rem:transition-to-hasse}
With the notion of lattice in hand, the Hasse diagram becomes a more powerful
visualisation tool: the Hasse diagram of a lattice has a unique bottom element
($\inf L$) and a unique top element ($\sup L$), and the join and meet of any
two elements can be read off as the unique path structure in the diagram.
The next section introduces Hasse diagrams, the duality principle, and the
connection between duality and the passage from $\sup$ to $\inf$.
\end{remark}